\section{Introduction}
\label{sec:intro}

The value of a text embedding model comes less from individual vectors than from the
geometry they induce: which examples are neighbours, and whether a similarity score means
the same thing for one query as for another. Since strong embedding models are increasingly
built on large language-model backbones \citep{chen2024bgem3,zhang2025qwen3embedding}, they
are expensive to serve, and distilling them into small encoders that preserve that geometry
is a practical necessity.

\paragraph{Existing methods differ in the target, not in the readout.}
The dominant approach is relational: rather than regressing student coordinates onto teacher
coordinates, the student is asked to preserve relations among examples
\citep{park2019rkd,tung2019sp}, in practice by matching a teacher-induced distribution over
candidate neighbours \citep{fang2021seed,tian2020crd,yang2020lsp}. Progress in this line has
come almost entirely from making that distribution carry more: distilling from several
teachers or several stages \citep{zhang2024jasperstella}, adding views and alternative
divergences \citep{xu2023distillcse,wen2023fdivergence}, supervising internal states under
capacity mismatch \citep{zhang2024dskd,dao2026talas}, or --- reading the teacher space as a
sampled manifold --- deriving targets at several resolutions from random walks on a
neighbourhood graph \citep{coifman2006diffusionmaps,klicpera2019gdc}. Whatever the target,
the student side is left untouched: a cosine profile over the candidate set is converted to
a distribution by \emph{one} softmax at \emph{one} temperature, and the several targets are
matched against that same distribution through a weighted sum of KL divergences. Our
starting point is that this uniform, unexamined choice is where these methods lose the
geometry they set out to transfer.

\paragraph{Three failures of the shared readout.}
None of the following is a statement about targets; each holds for any relational objective
of the form above. First, cross-entropy is linear in its target, so a weighted sum of KL
divergences against \emph{one} student distribution collapses onto the mixture of the
targets plus a target-only constant (Proposition~\ref{prop:collapse}). Enriching the target
family therefore buys smoothing and nothing else --- two families with the same mixture
induce identical gradients --- and the attained minimum is a generalised Jensen--Shannon
divergence among the targets, a number determined before any training
(Corollary~\ref{cor:floor}), so a loss curve that flattens is not evidence of convergence.
This applies verbatim to the multi-teacher, multi-view and multi-layer variants above.
Second, a softmax is invariant to a constant shift of its logits, so a single readout pins
the student's cosine profile only up to a per-anchor additive constant
(Proposition~\ref{prop:levels}): within-anchor ranking is transferred, absolute similarity
\emph{levels} are not. This is the precise reason relational distillation is observed to
preserve neighbour order while degrading tasks that apply one global threshold or rank all
pairs jointly, and it is what the usual remedy --- a pointwise term aligning $s_i$ to $t_i$
through a learned map --- provably cannot fix, since its invariance group is strictly larger
(Lemma~\ref{lem:gauge}). Third, targets derived from a sparsified graph assign mass zero
outside the anchor's support, where the cross-entropy gradient is strictly positive
(Proposition~\ref{prop:support}); the objective thus pushes apart precisely those pairs the
teacher scores as near but the mutual-neighbour test discarded --- so the standard defence
against hubness \citep{radovanovic2010hubs} creates a second, compounding failure.

\paragraph{One change addresses all three.}
Each failure is a statement about how the student reads a similarity vector out as a
distribution, so we intervene there. We replace the single softmax with a bank of softmaxes
at distinct temperatures, one per scale. Distinct temperatures break the algebraic identity
behind Proposition~\ref{prop:collapse}, and they over-determine the cosine profile, so the
minimiser depends on similarity levels rather than only on their shift-equivalence class
(Proposition~\ref{prop:levels}). For the support pathology we add no auxiliary term but
extend the scale family downwards: the zero-step member is the teacher's own dense kernel,
taken \emph{before} sparsification, and is the unique member whose support is the full
candidate set. The family thus interpolates between two extremes of one construction --- no
graph at $r=0$, deep walks at large $r$ --- rather than gluing a relational loss to a
diffusion loss. The result, \method, is a one-term objective with no balancing weight,
requiring only cached teacher embeddings: no hidden states, no online teacher pass, no
architecture change, no shared tokenizer.

Our contributions are the three results above --- a collapse identity with a pre-training
floor (Section~\ref{sec:collapse}), an under-determination result and its removal by
over-determination (Section~\ref{sec:levels}), and a support result with the zero-step
repair (Section~\ref{sec:support}) --- together with \method, the objective they force, and
build-time diagnostics that make each failure checkable before training
(Section~\ref{sec:method}). Because each proposition names the quantity left unconstrained,
it also says which downstream metrics should move when the corresponding component is
removed and which should not. We report results split along that line --- metrics defined by
a single global decision rule on cosine against metrics that use only within-anchor ordering
--- so the ablations test the analysis rather than merely rank configurations.
